\begin{proposition}[Albanese curves and parameter recovery]\hypertarget{tag-01Q3}{}\label{fam:src:A.4.1}\label{fam:st:A.4.1}
\StatementRecord{fam:src:A.4.1}{proved}{01Q3}{\SeriesRef{fam:src:A.1.2}, \SeriesRef{fam:src:A.2.1}}

With $\kappa_j(c)=\kappa_j^{\rm part}+c$, the Albanese elliptic curves of the reduced finite fibres are
\[
 \operatorname{Alb}(S_1)=\mathbb C/(\mathbb Z+\kappa_1(c)\mathbb Z),\qquad
 \operatorname{Alb}(S_2)=\mathbb C/(\mathbb Z+\tfrac12\kappa_2(c)\mathbb Z).
\]
Writing $J(q)=q^{-1}+744+196884q+\cdots$, their invariants are
\[
 j_1(w)=J(e^{-2\pi i\kappa_1^{\rm part}}w^2),\qquad
 j_2(w)=J(e^{-\pi i\kappa_2^{\rm part}}w).
\]
These are convergent formulas because the displayed nome moduli are below one. Outside the explicitly defined countable exceptional set in the proof, equality of both Albanese invariants forces $c'-c\in2\mathbb Z$. Thus the generic Albanese invariants recover the parameter modulo its exact period subgroup.
\end{proposition}


\begin{proof}
\begin{enumerate}[label=\arabic*.,leftmargin=*]
\item In the product covers of \ref{fam:src:A.1.2} the invariant holomorphic one-forms are exactly the multiples of $d\xi$: the moving factor is multiplied by a primitive third or fourth root. Integration thus identifies the Albanese with the quotient of the fixed line by the fixed-coordinate projections of every deck translation. These are, respectively,
\[
 \mathbb Z+\kappa_1\mathbb Z+\tfrac23\mathbb Z+\tfrac{\kappa_1}3\mathbb Z
   =\tfrac13\mathbb Z+\tfrac{\kappa_1}3\mathbb Z,
\quad
 \mathbb Z+\kappa_2\mathbb Z+\tfrac12\mathbb Z+\tfrac{\kappa_2}4\mathbb Z
   =\tfrac12\mathbb Z+\tfrac{\kappa_2}4\mathbb Z.
\]
Multiplication by three and two gives the stated maps of complex tori. The two displayed lattices include both the product-isogeny kernels and the affine quotient translations. Using the positive-imaginary representatives $-\kappa_1$ and $-\kappa_2/2$ in the standard modular $j$-series gives the formulas. At sufficiently small nome $J$ is injective: $1/J(q)=q+O(q^2)$ has nonzero derivative at zero. Thus $j_2$ alone distinguishes parameters within a sufficiently small punctured subdisc when both parameters lie there. For the full generic assertion put $x=-\kappa_2(c)/2$ and $d_*=\kappa_1^{\rm part}-\kappa_2^{\rm part}$, so $-\kappa_1(c)=2x-d_*$. The definitions $\kappa_1^{\rm part}=\beta^{\rm part}(z_1)+2\mu_1$ and
$\kappa_2^{\rm part}=\beta^{\rm part}(z_2)+3\mu_2$, together with
\cref{ari:C.1.2}, give
\[
 \operatorname{Im}d_*=
 -\frac{3}{2\pi}\log48+\frac3\pi\log6
 =-\frac{3}{2\pi}\log\frac43<0.
\] Equality of the two elliptic invariants at $c,c'$ gives matrices $M,N\in\operatorname{PSL}_2(\mathbb Z)$ with $x'=M(x)$ and
\[
 N(2x-d_*)=2M(x)-d_*.
\]
Define the exceptional set to be the $c\in\mathcal P$ obtained from roots of these equations for all pairs for which the two rational functions are not identical. Each equation has at most two roots after clearing denominators, so this is a countable, explicitly specified set. For a pair giving an identity, let $C(x)=2x-d_*$. Then $N=CMC^{-1}$. The real matrix $M$ preserves both the real circle $\mathbb R\cup\{\infty\}$ and $C^{-1}(\mathbb R\cup\{\infty\})$, the parallel line with imaginary part $\operatorname{Im}d_*/2$ together with infinity. Their unique intersection is infinity, so $M$ fixes infinity. An element of $\operatorname{PSL}_2(\mathbb Z)$ fixing infinity is $x\mapsto x+b$, $b\in\mathbb Z$, and then $N(y)=y+2b$. It follows that $c'-c=-2b$. This proves the necessity outside the specified exceptional set; \ref{fam:src:A.2.1} proves the necessity also on that set.
\end{enumerate}
\end{proof}